\documentclass[11pt]{article}

\usepackage[margin=1in]{geometry}
\usepackage[T1]{fontenc}
\usepackage[utf8]{inputenc}
\usepackage{array}
\usepackage{graphicx}
\usepackage{amsmath,amssymb}
\usepackage{hyperref}

\urlstyle{tt}
\newcommand{\tablefont}{\footnotesize}

\title{MindNav Experiment Record}
\author{}
\date{}

\begin{document}
\maketitle

\section{Purpose}

This document records the experimental path behind the MindNav paper. It is an
internal reference, not manuscript text. The goal is to make clear what was run,
why the next experiment was introduced, and what each stage showed.

The experiments proceeded in three stages:
\begin{enumerate}
    \item HM3D \textsc{val\_mini} 30-episode sanity checks.
    \item The original MindNav 50-episode evaluation on HM3D \textsc{val}.
    \item KG-complexity analysis, KG fixes, and a 50-episode heuristic control.
\end{enumerate}

Unless otherwise stated, SR denotes success rate and SPL denotes
success-weighted path length. All multi-robot runs use two robots.
Qwen2.5-7B denotes Qwen2.5-7B-Instruct deployed locally. All DeepSeek runs use
DeepSeek-V4-Flash through an API; there is no local DeepSeek deployment in
these experiments.

\section{Stage 1: HM3D \textsc{val\_mini} 30-Episode Sanity Checks}

The first stage was a small-split sanity check. The purpose was to verify that
the KG tool-calling pipeline could run end-to-end, compare it with Co-NavGPT
under similar infrastructure, and check whether the LLM backend made a visible
difference on a 30-episode split.

\begin{table}[t]
\centering
\tablefont
\caption{Stage 1: HM3D \textsc{val\_mini} sanity checks.}
\label{tab:stage1-val-mini}
\begin{tabular}{llrrrr}
\hline
Method & Backend / variant & Episodes & Successes & SR & SPL \\
\hline
MindNav KG tool-calling & Qwen2.5-7B & 30 & 22 & 0.733 & 0.344 \\
MindNav KG tool-calling & DeepSeek-V4-Flash API & 30 & 23 & 0.767 & 0.379 \\
Co-NavGPT & Qwen2.5-7B & 30 & 20 & 0.667 & 0.332 \\
Co-NavGPT & DeepSeek-V4-Flash API & 30 & 19 & 0.633 & 0.317 \\
MCoCoNav & Qwen2.5-VL-7B & 30 & 21 & 0.700 & 0.396 \\
\hline
\end{tabular}
\end{table}

At this stage, MindNav KG tool-calling was promising: it outperformed Co-NavGPT
on both Qwen and DeepSeek-V4-Flash API backends, and DeepSeek-V4-Flash was
slightly better than Qwen on this small split. However, the split is too small
to support a strong model
comparison. These runs mainly established that the KG pipeline was worth
evaluating on a larger fixed episode set.

\section{Stage 2: Original MindNav 50-Episode Evaluation}

The second stage evaluated MindNav KG and Co-NavGPT on a matched 50-episode
HM3D \textsc{val} matrix. This was the first more reliable comparison between
the KG tool-calling interface and the flat-text Co-NavGPT interface.

\begin{table}[t]
\centering
\tablefont
\caption{Stage 2: Original 50-episode HM3D \textsc{val} matrix.}
\label{tab:stage2-val50}
\begin{tabular}{llrrrr}
\hline
Method & Backend & Episodes & Successes & SR & SPL \\
\hline
MindNav KG & Qwen2.5-7B & 50 & 35 & 0.700 & 0.444 \\
MindNav KG & DeepSeek-V4-Flash API & 50 & 35 & 0.700 & 0.435 \\
Co-NavGPT & Qwen2.5-7B & 50 & 34 & 0.680 & 0.372 \\
Co-NavGPT & DeepSeek-V4-Flash API & 50 & 33 & 0.660 & 0.429 \\
\hline
\end{tabular}
\end{table}

MindNav KG improved over Co-NavGPT under both matched backends. With Qwen, the
gain was \(+0.020\) SR and \(+0.072\) SPL. With DeepSeek-V4-Flash API, the gain
was \(+0.040\) SR and \(+0.006\) SPL. The strongest SPL in this matrix was
MindNav KG with Qwen.

The Qwen--DeepSeek comparison was less straightforward. Both MindNav KG
backends achieved the same SR of 0.700, and Qwen had slightly higher SPL. At
the episode level, 29 episodes succeeded under both models, 9 failed under both
models, 6 succeeded only under Qwen, and 6 succeeded only under DeepSeek. Thus,
the aggregate success rate was tied, but the two models had different decision
behavior.

The working interpretation after this stage was that raw language-model
strength was not the only limiting factor. The LLM interface was narrow: the
model received active frontiers, room-type hints, nearby objects, target priors,
serialized KG text, recent history, and tool results, and it only had to assign
robots to frontiers. In this setting, stronger general reasoning may not
directly translate into higher SR/SPL if the final task is mostly stable
ranking and assignment under noisy perception.

\section{Stage 3: KG Complexity, KG Fixes, and Heuristic Control}

The third stage was motivated by the Stage 2 observation that DeepSeek did not
clearly outperform Qwen. We therefore inspected whether the KG itself was too
complex or noisy, changed the KG update procedure, and added a deterministic
heuristic policy as a control. The goal was not to introduce a new final method,
but to understand whether performance was limited by KG construction, LLM
frontier choice, or low-level navigation dynamics.

\subsection{Why We Suspected KG Complexity and Noise}

The earlier KG could mislead frontier assignment in several ways:
\begin{itemize}
    \item \textbf{Stale frontier state.} Frontier indices are step-local, but
    older room nodes could retain previous frontier metadata.
    \item \textbf{Stale inferred topology.} Edges such as
    \texttt{connected\_to}, \texttt{connected\_via}, \texttt{next\_to}, and
    \texttt{near} could persist after the map changed.
    \item \textbf{Object duplication.} The same object category could be
    inserted multiple times from frontier priors and semantic detections.
    \item \textbf{Weak confidence calibration.} Some object nodes used default
    confidence rather than confidence estimated from the semantic-map patch.
    \item \textbf{Weak wall/connectivity consistency.} Nearby room-grid cells
    could be connected even when wall evidence suggested separation.
    \item \textbf{Graph clutter.} Auxiliary object-to-room-type hints could add
    many extra edges and make the serialized KG harder to inspect.
\end{itemize}

These issues are especially relevant for a model that reads the KG structure
more aggressively. A model that mostly follows target priors may be insensitive
to KG noise; a model that tries to use topology, explored state, path hints, and
object--room relations can be helped by a clean KG but hurt by a noisy one.

\subsection{What We Changed in the KG}

The KG update was modified to make the graph more current and conservative:
\begin{itemize}
    \item clear \texttt{active\_frontier} markers at each global step and only
    mark currently extracted frontiers as active;
    \item rebuild inferred room topology from current evidence instead of
    accumulating stale room-connectivity edges;
    \item remove direct room connections that conflict with wall-separation
    evidence;
    \item estimate object confidence from the semantic-map patch and cap it at
    0.99;
    \item merge repeated same-category object detections with noisy-OR
    certainty,
    \[
        c_{\mathrm{new}} = 1 - (1-c_{\mathrm{old}})(1-c_{\mathrm{det}}),
    \]
    and deduplicate nearby object nodes after graph construction;
    \item store room-type suggestions as object properties rather than
    expanding them into many auxiliary edges;
    \item add a direct \texttt{TARGET\_<category>} node when the target category
    is visible with sufficient semantic-map evidence.
\end{itemize}

The intended effect was a cleaner decision state: fewer duplicate objects, fewer
stale inferred edges, more conservative topology, and active frontier ids that
matched the current planning step.

\subsection{Heuristic Policy Added as a Control}

We also added a deterministic KG heuristic. It uses the same optimized KG,
frontier enrichment, target-prior signal, and target shortcut as MindNav, but it
does not call an LLM. This made it possible to ask whether an explicit scoring
rule could match or outperform LLM frontier assignment.

If the target is directly represented by a \texttt{TARGET\_<target>} node, or if
an object of the target category has certainty above \(\tau=0.5\), the heuristic
sends the nearest robot to the frontier closest to the target. Otherwise, each
robot--frontier pair is scored as
\[
\begin{split}
s(r,f) ={}& 4.00\,p_{\mathrm{target}}(f)
          + 0.25\,\min(c_{\mathrm{room}}(f),1)
          + 0.15\,\min(A(f)/80,1) \\
        & - 0.35\,\min(d(r,f)/500,1)
          - 0.45\,n_{\mathrm{recent}}(f),
\end{split}
\]
where \(p_{\mathrm{target}}\) is the frontier target prior,
\(c_{\mathrm{room}}\) is room confidence, \(A\) is frontier area,
\(d(r,f)\) is robot--frontier distance, and \(n_{\mathrm{recent}}\) counts
recent repeated selections. For two robots, all distinct frontier pairs are
enumerated, with an additional diversity bonus for larger inter-frontier
distance.

\subsection{Stage 3 Results}

\begin{table}[t]
\centering
\tablefont
\caption{Stage 3: 50-episode reruns using the optimized MindNav KG pipeline.}
\label{tab:stage3-reruns}
\begin{tabular}{llrrrr}
\hline
Pipeline & Navigation brain / policy & Episodes & Successes & SR & SPL \\
\hline
MindNav optimized KG & Qwen2.5-7B LLM & 50 & 35 & 0.700 & 0.444 \\
MindNav optimized KG & DeepSeek-V4-Flash API & 50 & 32 & 0.640 & 0.384 \\
MindNav optimized KG & Deterministic KG heuristic & 50 & 36 & 0.720 & 0.422 \\
\hline
\end{tabular}
\end{table}

All three Stage 3 runs used the optimized MindNav KG pipeline. They differed
only in the navigation brain: Qwen2.5-7B, DeepSeek-V4-Flash API, or the
deterministic KG heuristic. With this optimized KG pipeline, Qwen reproduced the
Stage 2 Qwen result exactly, with \(\mathrm{SR}=0.700\) and
\(\mathrm{SPL}=0.444\). DeepSeek-V4-Flash dropped to \(\mathrm{SR}=0.640\) and
\(\mathrm{SPL}=0.384\). The heuristic
policy achieved the highest SR among the Stage 3 runs, \(\mathrm{SR}=0.720\),
although its SPL was lower than Qwen's SPL.

The main lesson from Stage 3 is that this setup is not a pure ``larger LLM is
better'' benchmark. A conservative policy that closely follows target prior,
room evidence, distance, recency, and frontier diversity is very competitive.
Qwen behaved closer to this conservative pattern, whereas DeepSeek appeared
more willing to deviate from the highest target-prior frontier after the KG
changes.

\section{Cross-Stage Interpretation}

Across the three stages, the experimental story is:
\begin{itemize}
    \item \textbf{Stage 1} showed that KG tool-calling was viable on
    \textsc{val\_mini} and generally better than Co-NavGPT on the same small
    split.
    \item \textbf{Stage 2} showed that MindNav KG improved over Co-NavGPT on a
    50-episode matrix, but also showed that DeepSeek did not clearly outperform
    Qwen2.5-7B.
    \item \textbf{Stage 3} showed that after KG cleanup, Qwen remained stable,
    DeepSeek became worse, and a deterministic heuristic achieved the best SR.
\end{itemize}

The likely bottleneck is therefore the embodied frontier-control interface, not
only language-model capacity. The LLM receives an already compressed state and
outputs a small assignment decision. Low-level navigation, perception quality,
target-prior calibration, KG noise, and post-processing all affect final SR/SPL.
In this regime, a model or policy that reliably follows the strongest numeric
evidence can outperform a model that explores more nuanced but noisy KG
structure.

\section{Detailed Decision Diagnostics}

The trace statistics support the Stage 3 interpretation. DeepSeek did not fail
because it ignored tools; it made more tool calls overall, partly because its
episodes lasted longer. The key behavioral difference was frontier selection.

\begin{table}[t]
\centering
\tablefont
\caption{Decision behavior on the Stage 3 val50 runs.}
\label{tab:decision-diagnostics}
\begin{tabular}{lrrrrrrr}
\hline
Variant & Steps & Target & Target rate & Same F & Top prior & Avg chosen & Avg max \\
\hline
KG heuristic & 386 & 76 & 0.197 & 0.039 & 0.969 & 0.396 & 0.483 \\
Qwen2.5-7B LLM & 352 & 67 & 0.190 & 0.028 & 0.972 & 0.427 & 0.527 \\
DeepSeek-V4-Flash API & 409 & 54 & 0.132 & 0.029 & 0.895 & 0.388 & 0.498 \\
\hline
\end{tabular}
\end{table}

\begin{table}[t]
\centering
\tablefont
\caption{KG tool-call counts on the Stage 3 LLM runs.}
\label{tab:tool-call-diagnostics}
\begin{tabular}{lrrrr}
\hline
Variant & Query objects & Estimate prob. & Assign & Check target \\
\hline
Qwen2.5-7B LLM & 782 & 782 & 284 & 67 \\
DeepSeek-V4-Flash API & 1028 & 1028 & 354 & 54 \\
\hline
\end{tabular}
\end{table}

Qwen chose a frontier with the highest target prior in 97.2\% of decision
steps, close to the heuristic's 96.9\%. DeepSeek did so in only 89.5\% of
decision steps. DeepSeek also had fewer target-found shortcut steps, more total
decision steps, and a lower average chosen prior. Same-frontier assignment rates
were low for all variants, so duplicate assignment of both robots to the same
frontier was not the primary failure mode.

These diagnostics suggest a concrete next direction: let the LLM estimate room
or frontier probabilities, but make the final frontier assignment deterministic,
with explicit distance, diversity, recency, and validity terms.

\end{document}
